\section{Техническая реализация программных модулей}

В этой главе мы разберем исходный код инструментов, которые используются для создания карт и навигации.

\subsection{Генератор карт (Marker Map Builder)}
Скрипт \texttt{marker\_map\_builder.py} позволяет создавать цифровые карты полигонов и экспортировать их в форматы JSON, XML и KML.

\subsubsection{Структура класса Marker}
Каждый маркер описывается классом \texttt{Marker}, который хранит:
\begin{itemize}
    \item \texttt{id}: Уникальный номер маркера.
    \item \texttt{size}: Физический размер стороны в метрах.
    \item \texttt{x, y, z}: Локальные координаты относительно начала координат карты.
    \item \texttt{yaw}: Угол поворота (ориентация).
    \item \texttt{dictionary}: Тип словаря (например, DICT\_6X6\_250).
\end{itemize}

\begin{codefiletitle}[python]{Класс Marker}{python/examples/advanced_navigation/marker_class_code.py}
\end{codefiletitle}

\subsubsection{Экспорт в XML}
Метод \texttt{save\_xml} генерирует структуру, понятную большинству промышленных систем. Мы используем f-строки для формирования XML-тегов.

\begin{codefiletitle}[python]{Экспорт в XML}{python/examples/advanced_navigation/save_xml_code.py}
\end{codefiletitle}

\subsubsection{Экспорт в KML для Google Earth}
Для отображения маркеров на глобусе Google Earth используется формат KML. Главная сложность здесь — перевод локальных метров $(x, y)$ в глобальные градусы широты и долготы.

Мы используем приближенную формулу (для малых расстояний):
\begin{itemize}
    \item 1 градус широты $\approx 111132$ метра.
    \item 1 градус долготы $\approx 111132 \cdot \cos(\text{широта})$ метров.
\end{itemize}

\begin{warnbox}{Точность GPS}
Данная формула имеет погрешность. Для профессиональной геодезии следует использовать библиотеку \texttt{pyproj} и проекцию UTM.
\end{warnbox}

\subsection{Модуль визуализации}
Для проверки правильности карты полезно отрисовать её на 2D-плоскости перед полетом. Библиотека \texttt{matplotlib} позволяет построить схему расположения маркеров.

\begin{codefiletitle}[python]{Визуализация карты}{python/examples/advanced_navigation/plot_map_code.py}
\end{codefiletitle}

\begin{taskbox}[Задание 7: Свой парсер]
Напишите функцию \texttt{load\_json(filename)}, которая считывает файл карты и создает объекты класса \texttt{Marker}. Проверьте, что карта корректно загружается после сохранения.
\end{taskbox}
